% orddiffops.tex

\begin{frame}
Ordinary Differential Operators: Vectors and Functions

\begin{itemize}
  \item Vectors are functions and functions are vectors
  
\item A vector $\begin{bmatrix}x_1\\x_2\\x_3\end{bmatrix}$ in 3-space is really a 
function $x:\set{1,2,3}\rightarrow \Real$
  
\item We just write them differently than functions, i.e., instead of $x(i)$ 
we write $x_i$. 

\end{itemize} 

\end{frame}


\begin{frame}
Ordinary Differential Operators: Vectors and Functions
 
\begin{itemize}
  
\item And since the domain $\set{1,2,3}$ of the function $x$ is so small, we can 
exhibit all values of $x$ explicitly in one fell swoop, i.e., 
as $\begin{bmatrix}x_1\\x_2\\x_3\end{bmatrix}$
  
\item But, by the same token, ordinary real functions $f:\Real\rightarrow\Real$ are 
vectors
  
\item They just have so many components that we can't exhibit them all at once

\end{itemize} 

\end{frame}


\begin{frame}
Ordinary Differential Operators: Vectors and Functions

\begin{itemize}
  
\item We can add vectors by adding their components: $z=x+y$ if and only if 
$z_i=x_i+y_i$ for each index $i$
  
\item By the same token, we can add functions componentwise:
  $h=f+g$ if and only 
if $h(t)=f(t)+g(t)$ for all $t$ in the domain
  
\item Componentwise multiplication of functions can also be defined 
(although this does not come up as often for ordinary vectors):  $h=f\cdot g$ if 
and only if $h(t)=f(t)\cdot g(t)$ for all $t$ in the domain

\end{itemize} 

\end{frame}



\begin{frame}
Ordinary Differential Operators: Vectors and Functions

\begin{itemize}
  
\item Finally, we can multiply functions by scalars, just like we do for vectors:  
$g=a\cdot f$ if and only if $g(t)=a\cdot f(t)$ for all $t\in\Real$.
  
\item The upshot of all this is that the set of all real functions is a vector space!
  
\item But we are really more interested in a subspace of this vector space, i.e., we really want to 
restrict our attention to the set of all \textbf{differentiable} functions on some closed interval.

\end{itemize} 

\end{frame}


\begin{frame}
Ordinary Differential Operators: Differentiable Functions

\begin{itemize}
  
\item Let $I=[a,b]$ be any closed interval of the real line
  
\item Consider the set $\Real^I=\set{f:I\rightarrow \Real}$ of all functions mapping $I=[a,b]$ into the real numbers
  
\item Then, as we have just pointed out, $\Real^I$ is a vector space over the reals

\end{itemize} 

\end{frame}



\begin{frame}
Ordinary Differential Operators
 
\begin{itemize}
  
\item We define the \textbf{\textit{derivative}} of a function $f$ at the point $t\in [a,b]$ as the 
limit as $s\rightarrow
 t$ of the difference quotient,
 \[\frac{f(s)-f(t)}{s-t},\]
assuming this limit exists.
  
\item Note that not all functions have derivatives!
  
\item Those that do have a derivative defined at every point of $I$ we call \textbf{\textit{differentiable}} on $I$.

\end{itemize} 

\end{frame}



\begin{frame}
Ordinary Differential Operators
 
\begin{itemize}
  
\item This operation turns out to be a linear operator on the vector space of functions that are differentiable 
on $I=[a,b]$.
  
\item In fact, we want to consider those functions that can be differentiated as many times as we like
  
\item Thus, we consider the space of infinitely differentiable, real-valued functions defined on the 
closed real interval $[a,b]$: $C^\infty[a,b]$
\end{itemize}

\end{frame}


% orddiffeqs.tex

\begin{frame}
Ordinary Differential Operators
 
\begin{itemize}
 
\item The operation of taking the derivative will enable us to develop deterministic differential
 equations.
 
\item Later we will define derivatives of stochastic processes in order to introduce and solve stochastic 
differential equations.  You cannot always take the derivative (or even integrate) a random process -- 
and sometimes not even the sample functions of a random process -- in the ordinary sense.

\end{itemize}

\end{frame}


% orddiffeqs.tex

\begin{frame}
Ordinary Differential Operators
 
\begin{itemize}
 
\item Since the derivatives of all orders are assumed to exist
 for all functions $f\in C^\infty[a,b]$, the 
limit of the
 difference quotient always exists for such functions.
 
\item Thus, we assign to each $f$ a new function $Df$ given by

\begin{equation}\label{deriv}
  [Df](t) = \lim_{s\rightarrow t}\frac{f(s)-f(t)}{s-t}
\end{equation}
 
\item This assignment defines a transformation
 \[D:C^\infty[a,b]\rightarrow C^\infty[a,b]\]
  that is a linear operator.

\end{itemize}

\end{frame}


\begin{frame}
Ordinary Differential Operators
 
\begin{itemize}
    
\item $D$ is linear: $D[af+bg]=a[Df] + b[Dg]$
    
\item Sometimes we get lax and simply write $Df(t)$ instead of $[Df](t)$.  However, keep in mind that $Df(t)$ 
says "Differentiate $f$ and then evaluate the resulting function at the point $t$."
    
\item $D$ obeys all the usual rules of differential calculus such as the product rule and the chain rule:

\[D(fg)=[Df]g+f[Dg]\]
\[D(f\circ g) = [Df]\circ g \cdot Dg\]
 
\end{itemize}

\end{frame}


\begin{frame}
Ordinary Differential Operators: Examples
 
\begin{itemize}
  
\item If $\id$ is the identity function, i.e., $\id(t)=t$, then
 $D\id=1$. (You probably know this 
as $\frac{dt}{dt}=1$; Note that 1 is the function whose value at every $t$ is 1.)
  
\item Using the product rule, we have, for example,

\begin{align*}
D\id^2 &= D\left[\id \cdot \id \right]\\
&= D\id \cdot \id + \id \cdot D\id\\
&= 1 \cdot \id + \id \cdot 1\\
&= \id + \id\\
&= 2\id
\end{align*}

(In more familiar notation this is just $\frac{dt^2}{dt}=2t$.)
 
\end{itemize}

\end{frame}


% orddiffeqs.tex

\begin{frame}
Ordinary Differential Operators
 
\begin{itemize}

 \item Note that we sometimes write $d/dt$ or $\frac{d}{dt}$ for
 the operator $D$ when we want to emphasize 
the independent
 variable $t$.  Or, instead of $\frac{dy}{dt}$ we write simply $\dot{y}$.
 
\item However, that is really unnecessary, adds no new information, and is just extra
 baggage.
 
\item Nevertheless, we shall do it sometimes because it may be more familiar.

\end{itemize}

\end{frame}


% orddiffeqs.tex

\begin{frame}
Ordinary Differential Operators
 
\begin{itemize}
 
\item Since there is really only one generating operator $D$ that we are interested in here,
 
we are less interested in multiplication than we are in function
 composition or iteration of this operator. 
For instance,
 $D^2f=[D\circ D](f) = D(Df)$.
 
\item By induction we have $D^n=\underbrace{D\circ\cdots \circ D}_n$.
 
\item Just as real-valued functions inherit algebraic operations from those of their range
 spaces, 
so too do differential operators:
 \[\left[aD^m+bD^n\right](f)=a(D^mf)+b(D^nf)\]

\end{itemize}

\end{frame}


% orddiffeqs.tex

\begin{frame}
Ordinary Differential Operators
 
\begin{itemize}
 
\item We can therefore generate polynomials in $D$:
 \[p(D)=a_nD^n+\cdots +a_1D+a_oI,\]
 
where $I$ is the identity operator on $C^\infty[a,b]$.  We
 usually just drop the identity operator, 
since it amounts to multiplying by 1.
 
\item Such polynomials in $D$ are still linear operators.
 
\item Thus, we have rediscovered the group
 ring $\Real(\Diff)$, where $\Diff$ in this case is the 
infinite cyclic
 group generated by the differential operator $D$.  We'll call
 this group ring 
 
the \textit{\textbf{polynomial ring of differential operators}}.

\end{itemize}

\end{frame}
